\documentclass[11pt]{article}
\usepackage[margin=1in]{geometry}
\usepackage{booktabs}
\usepackage{array}
\usepackage{hyperref}
\title{cot-loop PR\#2 Follow-up Audit \\ Last-token Final vs Mean-pool Final}
\author{}
\date{2026-03-03 06:52 UTC}

\begin{document}
\maketitle

\section*{Question Scope}
Follow-up asks addressed in this audit:
\begin{enumerate}
  \item Is ROC-AUC the right metric, and what do F1/accuracy say under class imbalance?
  \item Why does \texttt{last\_token\_final} show ID minus OOD ROC-AUC of \texttt{-0.1151}?
  \item Restrict conclusions to \texttt{last\_token\_final} and \texttt{mean\_pool\_final}.
  \item Can one rollout-built dataset be reused for both feature views?
\end{enumerate}

\section*{Observed Metrics (PR\#2 runs)}
\begin{center}
\begin{tabular}{@{}llrrrrr@{}}
\toprule
Feature & Split & Positives/60 & ROC-AUC & Accuracy & Macro-F1 & Maj. Acc. \\
\midrule
Last-token final & OOD & 19 & 0.6521 & 0.4833 & 0.4820 & 0.6833 \\
Last-token final & ID  & 6  & 0.5370 & 0.8000 & 0.5148 & 0.9000 \\
Mean-pool final  & OOD & 18 & 0.4603 & 0.5667 & 0.4841 & 0.7000 \\
Mean-pool final  & ID  & 5  & 0.7818 & 0.9167 & 0.4783 & 0.9167 \\
\bottomrule
\end{tabular}
\end{center}

\noindent
Majority-class accuracy is from always predicting non-loop (negative), computed from prevalence.

\section*{Interpretation}
\begin{itemize}
  \item \textbf{ROC-AUC should be kept, but not used alone.} It measures ranking quality and is threshold-free.
  \item \textbf{Accuracy is misleading on ID} because positive prevalence is only 5--6 out of 60. High ID accuracy can come from majority-negative predictions.
  \item \textbf{Macro-F1 gives thresholded class-balance signal} and remains around 0.48--0.51 across the four runs, indicating weak minority-class performance.
\end{itemize}

Recommended reporting bundle for this setup:
\begin{itemize}
  \item ROC-AUC
  \item Macro-F1
  \item Positive-class precision and recall
  \item Class prevalence
  \item PR-AUC (Average Precision) for imbalance-sensitive ranking
\end{itemize}

\section*{Why \texttt{last\_token\_final} has ID-\!OOD = -0.1151}
The measured values are OOD 0.6521 vs ID 0.5370, with only 6 ID positives. With this small positive count, AUC uncertainty is large and overlapping with OOD-scale uncertainty. So this negative gap is plausibly sampling variance, not strong evidence of an actual ID regression.

A second confound in PR\#2 is that each feature run rebuilt labels independently, so positives drifted across runs (OOD 18--21, ID 3--6), making feature comparisons directional rather than strictly matched.

\section*{Efficiency: one rollout build for both feature views}
Yes. The current dataset builder couples rollout-label generation and feature extraction, so changing feature view can trigger relabeling. A better design is:
\begin{enumerate}
  \item Build and cache one shared \textbf{split+label artifact} keyed by prompt source, split seed, rollout config, and loop definition.
  \item Build separate \textbf{feature-view shards} (last-token final, mean-pool final) against the same sample IDs and label order.
  \item Train probes by selecting feature key, without rerunning rollouts.
\end{enumerate}

Benefits:
\begin{itemize}
  \item One rollout job supports both feature views.
  \item No per-feature label drift.
  \item Cleaner apples-to-apples ID/OOD comparison.
\end{itemize}

\section*{Actionable Next Step}
For the next PR update, keep only \texttt{last\_token\_final} and \texttt{mean\_pool\_final}, add PR-AUC plus positive precision/recall reporting, and implement the shared-label-cache split so both views reuse one rollout build.

\end{document}
